\subsection{Структурная схема системы управления}

Для формирования структурной схемы системы управления выразим неизвестные величины из системы уравнений в операторной форме.

Из первого уравнения получим:
\begin{equation}
    \Phi_{\text{рпр}}
    =
    \frac{
        M_{\text{дпр}}
        +
        (bp+c)\Phi_{\text{м}}
    }{
        J_{\text{дпр}}p^2+bp+c
    }.
\end{equation}

Из второго уравнения:
\begin{equation}
    \Phi_{\text{м}}
    =
    \frac{
        bp+c
    }{
        J_{\text{м}}p^2+bp+c
    }
    \Phi_{\text{рпр}}.
\end{equation}

Из динамической характеристики двигателя:
\begin{equation}
    M_{\text{дпр}}
    =
    \frac{
        r_{\text{пр}}U
        -
        s_{\text{пр}}p\Phi_{\text{рпр}}
    }{
        \tau p+1
    }.
\end{equation}

Из уравнения системы управления:
\begin{equation}
    U
    =
    -kk_{\text{п}}\Phi_{\text{рпр}}
    -
    \frac{kk_{\text{и}}}{p}\Phi_{\text{м}}
    +
    \frac{kk_{\text{и}}}{p}\Phi^*.
\end{equation}

Полученные выражения позволяют построить структурную схему системы управления, представленную на рисунке~\ref{fig:pdf/structural_scheme}.

\screenshot[1.0]{pdf/structural_scheme}{Структурная схема системы управления}
